\documentclass[11pt,a4paper]{article}

% ============================================================
% Uniform MTDF preamble -- identical across all papers
% ============================================================
\usepackage[utf8]{inputenc}
\usepackage[T1]{fontenc}
\usepackage{lmodern}
\usepackage[english]{babel}
\usepackage[a4paper,margin=2.2cm]{geometry}
\usepackage{amsmath,amssymb,amsthm}
\usepackage{graphicx}
\usepackage{float}
\usepackage{booktabs}
\usepackage{array}
\usepackage{longtable}
\usepackage{tabularx}
\usepackage{xltabular}
\usepackage{multirow}
\usepackage{caption}
\usepackage[dvipsnames,table]{xcolor}
\usepackage{mdframed}
\usepackage{parskip}
\usepackage{ragged2e}
\usepackage[hidelinks,breaklinks=true]{hyperref}
\usepackage{microtype}
\usepackage{url}
% Allow \path{} to break more aggressively inside narrow table columns:
\makeatletter
\g@addto@macro\UrlBreaks{\do\_\do\-\do\.\do\/\do\:}
\makeatother

% Colours matching the HTML theme
\definecolor{accent}{HTML}{1A4B8A}
\definecolor{callout-bg}{HTML}{FBFBFB}
\definecolor{tablehead}{HTML}{F0F0F0}

% Prevent any table from overflowing
\setlength{\tabcolsep}{4pt}
\renewcommand{\arraystretch}{1.15}

% Caption style (compact, italic, small like the HTML)
\captionsetup{font=small,labelfont=bf,labelsep=period,justification=centerfirst}

% Hyperref metadata placeholders (filled per-paper below)
\hypersetup{
  pdftitle={Executive Briefing: Mesche's Tensor Dynamics Framework (MTDF)},
  pdfauthor={Ingo Mesche},
  pdfsubject={MTDF - Mesche Time-Dilation Field Theory},
  colorlinks=false
}

% Prevent overfull hbox from long URLs/DOIs in bibliography
\sloppy
\emergencystretch=3em

% Tolerant line breaking so long identifiers (Python filenames etc.) don't
% produce large overfull hboxes in narrow table columns. These settings keep
% all content on-page while slightly relaxing right-margin strictness.
\tolerance=1200
\hbadness=2000
\hyphenpenalty=150
\exhyphenpenalty=50

% ============================================================
\title{Executive Briefing: Mesche's Tensor Dynamics Framework (MTDF)}
\author{Ingo Mesche\\ \small Independent Researcher, Malta \\ \small \texttt{imesche@outlook.com}}
\date{V75 / July 2026}

\begin{document}
\maketitle
\begin{center}\itshape A reproducible, first-principles approach to the latent sector crisis\end{center}



\begin{mdframed}[linecolor=accent,linewidth=0.5pt,backgroundcolor=callout-bg,leftline=true,rightline=true,topline=true,bottomline=true,innerleftmargin=8pt,innerrightmargin=8pt,innertopmargin=6pt,innerbottommargin=6pt]
\textbf{Claim status.}
This paper distinguishes between:
(i) calibration anchors,
(ii) validation targets,
(iii) diagnostic consistency checks, and
(iv) speculative extensions.
Only results explicitly labelled as validation targets are used as evidence for the core MTDF claim. Diagnostic and speculative sections are included to define future tests, not to claim established confirmation.
\end{mdframed}

\section{The Premise}
For decades, observational anomalies in cosmology have been resolved by postulating invisible compositional components (Cold Dark Matter and Dark Energy). Mesche's Tensor Dynamics Framework (MTDF) proposes a dynamic alternative: spacetime acts as an elastic medium characterised by an intrinsic stress-energy tensor $\Sigma_{\mu\nu}$. By introducing just four structural parameters of differing epistemic status (see FAQ 4; including a strain-lever length scale $\beta = 22.7$ Mpc) plus the derived elastic modulus $E = (2/\alpha^2)\,\rho_c\,c^2 = 9.1 \times 10^{-10}$ Pa, MTDF recovers the cosmic expansion history and galactic dynamics without latent particle sectors.

\section{The Reproducible Evidence}
This framework is accompanied by a fully reproducible, GPU-accelerated computational pipeline encompassing modified Boltzmann solvers (class\_\allowbreak{}mtdf) and full MCMC chains. Key empirical results include:

\subsection{CMB Consistency}
MTDF remains consistent with the full Planck 2018 plik TTTEEE + low-$\ell$ + lensing likelihood as an external consistency check. Full MCMC runs against a matched $\Lambda$CDM control chain (identical binary, likelihoods, priors, and stop rules) yield $\Delta\chi^2 = -1.29$ at bounded best fits. This is a maximum-likelihood comparison, not a Bayes factor or Bayesian evidence calculation. The framework is consistent with early-universe physics: Planck 2018 does not distinguish MTDF from $\Lambda$CDM.

\subsection{The Dual-Epoch $H_0$ Mechanism}
MTDF predicts an early-universe energy injection of approximately $0.33\%$ (estimated geometrically as $\gamma_{\text{geom}} = 1/12$: one exact identity plus two order-unity heuristic factors; order-unity theoretical uncertainty), shifting the global CMB $H_0$ to $\approx 67.45$ km/\allowbreak{}s/\allowbreak{}Mpc. Simultaneously, local void stress depletion enhances local measurements; the result is consistent with SH0ES conditional on the depth of the local underdensity (a deep KBC-like void gives $H_{\text{local}} = 73.1 \pm 1.0$ km/\allowbreak{}s/\allowbreak{}Mpc, within $1\sigma$ of SH0ES; shallow reconstructions give $\approx 70$--$71$). MTDF does not predict the void itself.

The relevant classification is therefore not simply one $H_0$ value against another, but a method-dependent hierarchy: measurements calibrated through the local void environment are naturally shifted relative to global-background methods. Recent compilations of sound-horizon-free $H_0$ measurements are consistent with this pattern, with distance-ladder methods clustering near 73 km/\allowbreak{}s/\allowbreak{}Mpc while methods with less local-void exposure return lower values.

\subsection{Galactic Dynamics with a Frozen Parameter Set}
Using a frozen parameter set (no per-galaxy fitting), the framework reproduces the SPARC dataset rotation curves with an RMS log velocity scatter of 0.185 dex. In the tested galaxy-galaxy lensing implementations, it gives lower residuals than the corresponding $\Lambda$CDM NFW baseline profiles across KiDS and SDSS comparisons. These comparisons should be read as model-implementation tests, not as a final exclusion of the full $\Lambda$CDM modelling space.

\subsection{The $S_8$ Posture}
The derived $S_8$ is $0.8473 \pm 0.0138$ (native, sealed v2 posterior) against $0.830 \pm 0.013$ for $\Lambda$CDM: MTDF predicts a sharp, \emph{higher} clustering amplitude, in tension with $\Lambda$CDM-compressed weak-lensing inferences (KiDS-1000, DES-Y3); the sealed posture is S8-WORSENED. Whether those compressions represent the true amplitude in an MTDF universe is under preregistered adjudication (FWD-S8), and the late-time growth history is a registered, falsifiable prediction (GH-1).

\subsection{Solar System Safety}
MTDF's structural suppression keeps all Post-Newtonian (PPN) parameters safely below observational bounds: $\gamma - 1 \approx 2\times10^{-8}$, three orders of magnitude under the Cassini bound $2.3\times10^{-5}$ (sealed safety ledger). Wide binary tests are consistent with effectively Newtonian dynamics ($\gamma = 1.10 \pm 0.01$, below the 1.15 threshold; $\gamma$ tracks the contamination fraction (r = 0.99) and falls to 1.06--1.08 under quality cuts). This remains an external discriminator: MTDF and acceleration-threshold models (e.g. MOND) agree on galaxy-scale phenomenology but differ for wide binaries, where MTDF predicts an EFE-modulated transition (onset near $\sim 8.6$ kAU) rather than a MOND-like universal boost; the current observational situation is mutually contested between datasets, and the row is carried as PREDICTION-STATED / DATA-CONTESTED in the safety ledger.

\section{Falsifiable Prediction}
MTDF predicts that Type Ia Supernova distance residuals should exhibit an environment-linked transition at low redshift, sharply confined below $z \approx 0.04$, rather than a uniform global correlation; the conjectured link to the strain-lever length scale $\beta = 22.7$ Mpc remains a hypothesis (the transition scale itself is comoving \~{}170 Mpc). A merged analysis of 4,188 supernovae (Pantheon+, ZTF DR2, Foundation DR1) across three independent void catalogues (VoidFinder, REVOLVER, VIDE) is consistent with this: the global constant-$\gamma$ model is null, while a piecewise transition at $z \approx 0.04$ yields $\Delta\chi^2 = 36$--$39$ ($p < 0.001$, all three catalogues; 36.8 nominal, diagonal covariance). SALT2 population controls (stretch, colour) shift the result by less than 0.3$\sigma$, and the same physical supernovae drive the signal across all three void reconstructions.

\section{Data and code availability}
The complete MTDF reproducibility package, including the modified CLASS solver
(\texttt{class\_\allowbreak{}mtdf}), the V19 strict validation workbook, the
\texttt{{\small\protect\path{run_validate.py}}} pipeline, the gravity-sector artefact manifest,
and the Phase 2 CosmoPower emulator cache, is archived on Zenodo
(concept DOI \href{https://doi.org/10.5281/zenodo.19741058}{{\small\protect\path{10.5281/zenodo.19741058}}};
this release, \texttt{v1.1.7}, DOI
\href{https://doi.org/10.5281/zenodo.21525855}{{\small\protect\path{10.5281/zenodo.21525855}}})
and mirrored on GitHub at
\href{https://github.com/IMesche/MTDF}{github.com/\allowbreak{}IMesche/\allowbreak{}MTDF}.
The \texttt{v1.1.7} Git tag corresponds to the archived Zenodo release. All
numerical results reported here can be regenerated from the shipped workbook
and scripts following the procedure in the top-level \texttt{{\small\protect\path{README.md}}}.

All assertions in the accompanying DOI publications are backed by a fully open codebase. Researchers can download the data, run the MCMC chains, and test the MTDF constraints directly.

\begin{mdframed}[linecolor=accent,linewidth=0.5pt,backgroundcolor=callout-bg,leftline=true,rightline=true,topline=true,bottomline=true,innerleftmargin=8pt,innerrightmargin=8pt,innertopmargin=6pt,innerbottommargin=6pt]
\textbf{Repository \& Reproducibility:}\begin{itemize}
  \item \textbf{Source code:} \href{https://github.com/IMesche/MTDF}{github.com/\allowbreak{}IMesche/\allowbreak{}MTDF}
  \item \textbf{Papers:} See \texttt{papers/\allowbreak{}} directory (MTDF\_\allowbreak{}01 through MTDF\_\allowbreak{}08)
  \item \textbf{Reproduce:} \texttt{bash {\small\protect\path{setup_environment.sh}} \&\& bash {\small\protect\path{scripts/download_data.sh}} \&\& cd validation/\allowbreak{}code \&\& python {\small\protect\path{run_validate.py}}}
  \item \textbf{DOI:} \emph{See CITATION.cff in the repository root}
\end{itemize}
\end{mdframed}

\section{Frequently Asked Questions}

\subsection{Is MTDF a modified-gravity theory?}
MTDF is a field theory in which spacetime possesses an elastic stress-energy component. It does not modify the geodesic motion of particles or photons relative to General Relativity. Instead, it introduces an additional, physically specified stress contribution that alters the effective dynamics at large scales.

\subsection{Does MTDF eliminate dark matter entirely?}
MTDF does not require a particle dark matter component at late times and on galactic scales, where the field's environment-dependent dynamics reproduce the gravitational effects commonly attributed to dark matter as an emergent consequence of spacetime stress sourced by baryonic matter. In the early universe, the CMB fits retain an $\Omega_c$-like density parameter ($\Omega_c h^2 \approx 0.12$, the $\Lambda$CDM value), which in the current implementation is interpreted as the stress-field energy component constructed to be degenerate with cold dark matter at early times, not as a particle species. A first-principles derivation showing that the developed field relaxes into the late-time phenomenology from this early degenerate phase is an open theoretical task.

\subsection{How does MTDF reproduce the CMB power spectrum?}
MTDF is constructed to be effectively degenerate with $\Lambda$CDM at early times. The model reproduces the observed CMB power spectrum using a derived early-field energy fraction ($f_{\text{kick}} \approx 0.33\%$), without calibrating parameters on CMB data. CMB observations are treated as external validation, not input.

\subsection{What are the free parameters of MTDF?}
MTDF has four fundamental parameters: $\alpha$, $\beta$, $\tau$, and $\beta_{\text{eos}}$, and they do not share a single epistemic status. The strain-lever length scale $\beta$ is anchored to an independent observation (the void radius distribution); the coupling $\alpha$ is a frozen, void-motivated value that is not independently measured; the relaxation time $\tau$ is synchronised to the cosmic age; and the equation-of-state parameter $\beta_{\text{eos}}$ is a phenomenological fit. None is calibrated on CMB data. Additional quantities, including the elastic modulus $E$ and photon coupling $\kappa$, are derived or theory-motivated and are not free parameters.

\subsection{How does MTDF differ from MOND and other modified gravity models?}
MTDF differs in both mechanism and predictions. In particular: MTDF does not introduce an acceleration threshold; MTDF predicts an EFE-modulated wide-binary transition (onset near $\sim 8.6$ kAU) rather than a MOND-like universal boost (the observational record is contested; PREDICTION-STATED / DATA-CONTESTED); and MTDF produces environment-dependent effects tied to large-scale structure. These distinctions provide direct observational discriminators between MTDF and MOND-like frameworks.

\subsection{Does MTDF make new, testable predictions?}
Yes. The framework predicts: a correlation between supernova residuals and large-scale environment; a sharp transition in this signal at low redshift ($z \approx 0.04$); near-Newtonian wide-binary behaviour over most of the sampled range, with an EFE-modulated transition setting in near $\sim 8.6$ kAU (PREDICTION-STATED / DATA-CONTESTED); and a method-dependent hierarchy in $H_0$ measurements. These predictions are falsifiable with current or near-future datasets.

\subsection{Is MTDF consistent with Solar System tests?}
Yes. The predicted deviations from General Relativity at Solar System scales sit orders of magnitude below current experimental sensitivity (PPN $\gamma - 1 \approx 2\times10^{-8}$ against the Cassini bound $2.3\times10^{-5}$; SEP $\eta \approx 2\times10^{-8}$ against the LLR bound $4.4\times10^{-4}$).

\subsection{Does MTDF resolve the Hubble tension?}
MTDF provides a physical mechanism that produces a dual-epoch behaviour in the expansion rate. Early-universe observables remain consistent with CMB-inferred values, while late-time, environment-sensitive measurements naturally produce higher local $H_0$ values, conditional on the depth of the local underdensity; shallow reconstructions give roughly 70 to 71 km/\allowbreak{}s/\allowbreak{}Mpc. This leads to a method-dependent hierarchy that matches current observational patterns.

\subsection{What would falsify MTDF?}
MTDF would be falsified if: no statistically significant environment-dependent signal is observed in low-redshift supernova data; no transition is observed near $z \approx 0.04$ in sufficiently large samples; or the EFE-modulated wide-binary transition near $\sim 8.6$ kAU is confirmed absent at decisive precision; a sustained, confirmed MOND-like universal boost in wide binaries would equally contradict the prediction. These conditions are explicitly testable.

\subsection{What is the minimal claim of MTDF?}
MTDF proposes that a single, physically defined spacetime stress field can account for both cosmological expansion and galaxy-scale dynamics, while remaining consistent with current high-precision observations and making testable predictions at low redshift.

This document accompanies the full MTDF paper suite (MTDF\_\allowbreak{}01 through MTDF\_\allowbreak{}08). For the complete technical treatment, please refer to the published papers.


\end{document}
